% 第2章 孤立磁矩（Isolated magnetic moments）
\chapter{孤立磁矩}

本章考察孤立磁矩的性质。在此阶段，我们忽略不同原子的磁矩之间的相互作用，也忽略
磁矩与其近邻环境之间的相互作用，剩下的就只是孤立原子的物理，以及它们与外加磁场
的相互作用。这当然不意味着问题变简单了，但复杂性的来源是单个原子内电子的组合
方式，而不是凝聚态中原子数目巨大这一事实。经过这一简化，大量的原子只是使磁化率
这类性质多出一个因子 $n$——单位体积内的原子数。

\section{磁场中的原子}

1.1 节（式 1.35）已经证明：沿 $z$ 轴的磁场中，电子自旋的能量为
\begin{equation*}
E = g\mu_{\mathrm{B}}Bm_s\tag{2.1}
\end{equation*}
其中 $g\approx2$，$m_s=\pm\frac{1}{2}$，故 $E\approx\pm\mu_{\mathrm{B}}B$。除自旋
角动量外，原子中的电子还具有轨道角动量。若原子中第 $i$ 个电子的位置为 $\mathbf{r}_i$、
动量为 $\mathbf{p}_i$，则总角动量 $\hbar\mathbf{L}$ 为
\begin{equation*}
\hbar\mathbf{L} = \sum_i \mathbf{r}_i\times\mathbf{p}_i\tag{2.2}
\end{equation*}
求和遍及原子中所有电子。考虑哈密顿量为 $\hat{\mathcal{H}}_0$ 的原子：
\begin{equation*}
\hat{\mathcal{H}}_0 = \sum_{i=1}^{Z}\left(\frac{p_i^{2}}{2m} + V_i\right)\tag{2.3}
\end{equation*}
它是对原子中 $Z$ 个电子求和的电子动能（第 $i$ 个电子的 $p_i^{2}/2m_{\mathrm{e}}$）
与势能（第 $i$ 个电子的 $V_i$）之和。设 $\hat{\mathcal{H}}_0$ 的本征态与本征值
均为已知。

现在加入磁场
\begin{equation*}
\mathbf{B} = \nabla\times\mathbf{A}\tag{2.4}
\end{equation*}
其中 $\mathbf{A}$ 是磁矢势。我们选取规范\fn{1}{式 2.4 把 B 与 A 联系起来。但对给定的磁场 B，磁矢势 A 并非唯一确定：可以给 A 加上某个标量势的梯度而仍得到同样的 B。我们对 A 所作的这种选择称为规范选择。}使得
\begin{equation*}
\mathbf{A}(\mathbf{r}) = \frac{\mathbf{B}\times\mathbf{r}}{2}.\tag{2.5}
\end{equation*}
于是动能须按 1.2 节的方案修改。由于电子电荷为 $-e$，动能为
$[\mathbf{p}_i+e\mathbf{A}(\mathbf{r}_i)]^{2}/2m_{\mathrm{e}}$，受扰动的哈密顿量
应写成
\begin{align*}
\hat{\mathcal{H}} &= \sum_{i=1}^{Z}\left(\frac{[\mathbf{p}_i+e\mathbf{A}(\mathbf{r}_i)]^{2}}{2m_{\mathrm{e}}} + V_i\right) + g\mu_{\mathrm{B}}\mathbf{B}\cdot\mathbf{S}\tag{2.6}\\
&= \sum_i\left(\frac{p_i^{2}}{2m_{\mathrm{e}}} + V_i\right) + \mu_{\mathrm{B}}(\mathbf{L}+g\mathbf{S})\cdot\mathbf{B} + \frac{e^{2}}{8m_{\mathrm{e}}}\sum_i(\mathbf{B}\times\mathbf{r}_i)^{2}\tag{2.7}\\
&= \hat{\mathcal{H}}_0 + \mu_{\mathrm{B}}(\mathbf{L}+g\mathbf{S})\cdot\mathbf{B} + \frac{e^{2}}{8m_{\mathrm{e}}}\sum_i(\mathbf{B}\times\mathbf{r}_i)^{2}\tag{2.8}
\end{align*}
对原哈密顿量 $\hat{\mathcal{H}}_0$ 的主要微扰通常是 $\mu_{\mathrm{B}}(\mathbf{L}+g\mathbf{S})\cdot\mathbf{B}$
项，但正如将要看到的，它有时会消失。这是原子自身磁矩的效应，称为顺磁项（paramagnetic
term）。第三项 $(e^{2}/8m_{\mathrm{e}})\sum_i(\mathbf{B}\times\mathbf{r}_i)^{2}$ 来自
抗磁矩。这两项贡献将在 2.3 节（抗磁性）与 2.4 节（顺磁性）中详细讨论。下一节概述
需要解释的效应。

\section{磁化率}

如 1.1.4 节所示，对线性材料 $M=\chi H$，$M$ 是单位体积的磁矩（磁化强度），$\chi$
是磁化率（无量纲）。注意 $M$ 的定义意味着 $\chi$ 表示单位体积中由磁场 H 诱导的
磁矩。磁化率常以摩尔磁化率（molar magnetic susceptibility）$\chi_{\mathrm{m}}$ 列表：
\begin{equation*}
\chi_{\mathrm{m}} = \chi V_{\mathrm{m}}.\tag{2.9}
\end{equation*}
式中 $V_{\mathrm{m}}$ 是摩尔体积，即 1 摩尔（$6.022\times10^{23}$ 个化学式单位）
物质占据的体积。摩尔体积（单位 m$^3$）等于物质的相对原子质量\fn{2}{相对原子质量是 1 摩尔的质量。注意相对原子质量通常以克为单位列表。}（单位 kg）除以密度
$\rho$（单位 kg\,m$^{-3}$）。质量磁化率 $\chi_g$ 定义为
\begin{equation*}
\chi_g = \frac{\chi}{\rho},\tag{2.10}
\end{equation*}
单位为 m$^3$\,kg$^{-1}$。若干物质的磁化率数值列于表 2.1。磁化率为负者以抗磁性为主，
为正者以顺磁性为主。

\begin{table}
\centering
\caption{若干物质在 298 K 下的磁化率 $\chi$ 与摩尔磁化率 $\chi_{\mathrm{m}}$。水、苯和 NaCl
是弱抗磁体（磁化率为负）；$\mathrm{CuSO_4\cdot5H_2O}$、$\mathrm{MnSO_4\cdot4H_2O}$、Al
和 Na 是顺磁体（磁化率为正）。}
\begin{tabular}{lcc}
\toprule
 & $\chi/10^{-6}$ & $\chi_{\mathrm{m}}/10^{-10}$（m$^3$\,mol$^{-1}$）\\
\midrule
水 & $-90$ & $-16.0$\\
苯 & $-7.2$ & $-6.4$\\
NaCl & $-13.9$ & $-3.75$\\
石墨（$\parallel$） & $-260$ & $-31$\\
石墨（$\perp$） & $-3.8$ & $-4.6$\\
Cu & $-1.1$ & $-0.078$\\
Ag & $-2.4$ & $-0.25$\\
\midrule
$\mathrm{CuSO_4\cdot5H_2O}$ & $176$ & $192$\\
$\mathrm{MnSO_4\cdot4H_2O}$ & $2640$ & $2.79\times10^{3}$\\
Al & $22$ & $2.2$\\
Na & $7.3$ & $1.7$\\
\bottomrule
\end{tabular}
\end{table}

周期表前 60 个元素的磁化率绘于图 2.1。其中一些为负，表明 2.3 节将讨论的抗磁性
占主导；另一些为正，表明存在顺磁性，这将在 2.4 节讨论。

\begin{figure}
\centering
\includegraphics[width=0.75\textwidth]{fig2_01.pdf}
\caption{周期表前 60 个元素在室温下的质量磁化率，按原子序数排列。Fe、Co、Ni 是铁磁体，
在无外加磁场时也有自发磁化。}
\end{figure}

\section{抗磁性}

所有材料都表现出一定程度的抗磁性（diamagnetism）\fn{3}{前缀 dia- 意为“相对”“横跨”（英语的 diagonal（对角）、diameter（直径）等词即由此而来）。}——一种弱的、
负的磁化率。对抗磁物质，磁场诱导出的磁矩与产生它的外加磁场方向相反。

这一效应常被经典地讨论：磁场对电子轨道运动的作用产生反电动势\fn{4}{electromotive force（电动势）。}，由楞次定律，它反抗产生自己的磁场。然而，
上一章的玻尔--范列文定理提醒我们警惕这类论证——它们试图证明经典体系中磁场可以
诱导磁矩。\fn{5}{延伸阅读中有进一步讨论。}抗磁性是完全量子力学的现象，必须按量子
力学处理。

用量子力学方法很容易说明这一效应。考虑没有未满电子壳层的原子，此时式 2.8 中的
顺磁项可以忽略。若 $\mathbf{B}$ 平行于 $z$ 轴，则 $\mathbf{B}\times\mathbf{r}_i = B(-y_i, x_i, 0)$，且
\begin{equation*}
(\mathbf{B}\times\mathbf{r}_i)^{2} = B^{2}(x_i^{2}+y_i^{2})\tag{2.11}
\end{equation*}
于是抗磁项引起的基态能量一阶移动为
\begin{equation*}
\Delta E_0 = \frac{e^{2}B^{2}}{8m_{\mathrm{e}}}\sum_{i=1}^{Z}\langle 0|(x_i^{2}+y_i^{2})|0\rangle,\tag{2.12}
\end{equation*}
$|0\rangle$ 是基态波函数。若假定原子球对称\fn{6}{若总角动量 J 为零，这是一个良好的近似。}，则 $\langle x_i^{2}\rangle=\langle y_i^{2}\rangle=\frac{1}{3}\langle r_i^{2}\rangle$，可得
\begin{equation*}
\Delta E_0 = \frac{e^{2}B^{2}}{12m_{\mathrm{e}}}\sum_{i=1}^{Z}\langle 0|r_i^{2}|0\rangle.\tag{2.13}
\end{equation*}

考虑体积 $V$ 内由 $N$ 个离子（每个有 $Z$ 个质量为 $m$ 的电子）组成、所有壳层均填满
的固体。求 $T=0$ 时的磁化强度可以循附录 E 的做法：
\begin{equation*}
M = -\frac{\partial F}{\partial B} = -\frac{N}{V}\frac{\partial\Delta E_0}{\partial B} = -\frac{Ne^{2}B}{6m_{\mathrm{e}}}\sum_{i=1}^{Z}\langle r_i^{2}\rangle,\tag{2.14}
\end{equation*}
$F$ 是亥姆霍兹自由能。由此可提取抗磁磁化率 $\chi=M/H\approx\mu_0M/B$（设 $\chi\ll1$）：
\begin{equation*}
\chi = -\frac{N}{V}\frac{e^{2}\mu_0}{6m_{\mathrm{e}}}\sum_{i=1}^{Z}\langle r_i^{2}\rangle.\tag{2.15}
\end{equation*}
此式使用了一阶微扰论（二阶项将在 2.4.4 节考虑）。温度自零上升后，基态以上的态对
抗磁磁化率的贡献逐渐变大，但这是边缘效应。抗磁磁化率通常基本不依赖温度。

把实验测得的多种离子的抗磁摩尔磁化率对 $Z_{\mathrm{eff}}r^{2}$ 作图，可以相当粗糙地
检验上式，其中 $Z_{\mathrm{eff}}$ 是离子外壳层中的电子数\fn{7}{对离子而言这一数值不同于原子序数 Z，故用符号 $Z_{\mathrm{eff}}$ 表示“有效”原子序数；我们忽略了内壳层电子。}，$r$ 是实测离子半径。此近似假定离子外壳层的所有电子的
$\langle r_i\rangle^{2}$ 大致相同，即
\begin{equation*}
\sum_{i=1}^{Z_{\mathrm{eff}}}\langle r_i^{2}\rangle \approx Z_{\mathrm{eff}}r^{2}.\tag{2.16}
\end{equation*}

若干离子的抗磁磁化率示于图 2.2。实验值由比较一系列离子盐（NaF、NaCl、NaBr、KCl、
KBr、…）的抗磁磁化率推得。这一做法并不精确——离子中并非所有电子的平均半径平方都
相同（故式 2.16 绝非严格），但符合程度仍相当可观。之所以选离子，是因为例如 Na 和
Cl 原子有未成对电子，而 $\mathrm{Na^{+}}$ 与 $\mathrm{Cl^{-}}$ 离子都是闭壳层结构，
与 Ne、Ar 类似（可参照下文图 2.13 的周期表）。于是本会主导原子磁响应的顺磁效应
在离子中可以忽略。

\begin{figure}
\centering
\includegraphics[width=0.6\textwidth]{fig2_02.pdf}
\caption{若干离子的抗磁摩尔磁化率实验值 $\chi_m$ 对 $Z_{\mathrm{eff}}r^{2}$ 作图，
$Z_{\mathrm{eff}}$ 为离子外层电子数，$r$ 为实测离子半径。}
\end{figure}

在具有离域 $\pi$ 电子的分子（如萘和石墨）中，可以观察到较大且各向异性的抗磁磁化率。
萘由两个苯环沿一边并合而成（图 2.3(a)）。$\pi$ 电子非常活跃，诱导电流可以沿环边缘
流动，产生大的抗磁磁化率；当磁场垂直于环平面时该效应最大。

\begin{figure}
\centering
\includegraphics[width=0.45\textwidth]{fig2_03.pdf}
\caption{（a）萘由两个稠合的苯环组成。（b）石墨由六角层堆叠而成，碳原子画成黑点。
碳原子在隔层（而非相邻层）之间对齐（见竖直虚线）。}
\end{figure}

有效环直径是原子直径的数倍，故效应很大。石墨同理：它由结合疏松的六角层构成
（图 2.3(b)）。磁场垂直于层面时的抗磁磁化率远大于平行方向。

抗磁性存在于一切材料之中，但它是弱效应：要么可以忽略，要么只是对更强效应的小
修正。

\section{顺磁性}

顺磁性（paramagnetism）\fn{8}{前缀 para- 意为“伴随”“沿着”（英语的 parallel（平行）等词即由此而来）。}对应正的磁化率：外加磁场诱导的磁化强度与磁场平行。
上一节考虑的是不含未成对电子的材料，其原子或分子在未加场时没有磁矩。本节关心
的是因未成对电子而具有非零磁矩的原子。无外加磁场时，这些磁矩指向随机——相邻原子
的磁矩之间相互作用极弱，可视为彼此独立。加上磁场后磁矩被排齐，排齐的程度（因而
诱导的磁化强度）取决于磁场强度。

原子磁矩与原子的总角动量 $\mathbf{J}$ 相联系，$\mathbf{J}$ 是轨道角动量 $\mathbf{L}$
与自旋角动量 $\mathbf{S}$ 之和：
\begin{equation*}
\mathbf{J} = \mathbf{L} + \mathbf{S}.\tag{2.17}
\end{equation*}
与全书一致，这些量都以 $\hbar$ 为单位度量。自旋与轨道部分如何组合将在后面几节
详述；本节只假定每个原子具有大小为 $\mu$ 的磁矩。

增强磁场会使自旋排齐，升温则使它们随机化，因此可以预期顺磁材料的磁化强度依赖于
比值 $B/T$。顺磁效应一般远强于抗磁效应，尽管抗磁性总是作为弱的负贡献存在。

\subsection{顺磁性的半经典处理}

我们先半经典地处理顺磁性（下面会看到这对应 $J=\infty$），忽略量子化导致的“磁矩
只能指向某些方向”这一事实。考虑与外场 $\mathbf{B}$（不失一般性取沿 $z$ 方向）成
$\theta$ 到 $\theta+\mathrm{d}\theta$ 角的磁矩，其能量为 $-\mu B\cos\theta$，沿
$\mathbf{B}$ 的净磁矩为 $\mu\cos\theta$。若磁矩可以完全随机地选择方向，则取角介于
$\theta$ 与 $\theta+\mathrm{d}\theta$ 之间的比例正比于图 2.4 中环带的面积——单位
半径球面上为 $2\pi\sin\theta\,\mathrm{d}\theta$。单位球总面积为 $4\pi$，故该比例
为 $\frac{1}{2}\sin\theta\,\mathrm{d}\theta$。温度 $T$ 下磁矩取角介于 $\theta$ 与
$\theta+\mathrm{d}\theta$ 的概率正比于统计因子 $\frac{1}{2}\sin\theta\,\mathrm{d}\theta$
与玻尔兹曼因子 $\exp(\mu B\cos\theta/k_{\mathrm{B}}T)$ 之积（$k_{\mathrm{B}}$ 为玻尔
兹曼常数）。于是沿 $\mathbf{B}$ 的平均磁矩为
\bio{Paul Langevin\\（1872--1946）}
\begin{align*}
\langle\mu_z\rangle &= \frac{\int_0^{\pi}\mu\cos\theta\,\exp(\mu B\cos\theta/k_{\mathrm{B}}T)\,\frac{1}{2}\sin\theta\,\mathrm{d}\theta}{\int_0^{\pi}\exp(\mu B\cos\theta/k_{\mathrm{B}}T)\,\frac{1}{2}\sin\theta\,\mathrm{d}\theta}\tag{2.18}\\
&= \mu\,\frac{\int_{-1}^{1}x e^{yx}\,\mathrm{d}x}{\int_{-1}^{1}e^{yx}\,\mathrm{d}x},\tag{2.19}
\end{align*}
其中定义 $y=\mu B/k_{\mathrm{B}}T$，$x=\cos\theta$。由此得
\begin{equation*}
\frac{\langle\mu_z\rangle}{\mu} = \coth y - \frac{1}{y} \equiv L(y)\tag{2.20}
\end{equation*}
$L(y)=\coth y-1/y$ 称为朗之万函数（Langevin function），绘于图 2.5。$y$ 小时
\begin{equation*}
\coth(y) = \frac{1}{y} + \frac{y}{3} + O(y^{3})\tag{2.21}
\end{equation*}
于是
\begin{equation*}
L(y) = \frac{y}{3} + O(y^{3}).\tag{2.22}
\end{equation*}

\begin{figure}
\centering
\includegraphics[width=0.42\textwidth]{fig2_04.pdf}
\caption{计算顺磁材料的平均磁矩：考虑磁矩与 $z$ 轴夹角介于 $\theta$ 与
$\theta+\mathrm{d}\theta$ 的概率，它正比于单位球上阴影环带的面积
$2\pi\sin\theta\,\mathrm{d}\theta$。}
\end{figure}

\begin{figure}
\centering
\includegraphics[width=0.5\textwidth]{fig2_05.pdf}
\caption{经典顺磁体的磁化强度由朗之万函数 $L(y)=\coth y-\frac{1}{y}$ 描述。$y$ 小时
$L(y)\approx y/3$，如图中在原点附近与曲线相切的直线所示。磁场增强或温度降低时，
磁化强度的数值增大。}
\end{figure}

用 $n$ 表示单位体积的磁矩数。饱和磁化强度 $M_{\mathrm{s}}$ 是所有磁矩完全排齐时
能得到的最大磁化强度，$M_{\mathrm{s}}=n\mu$。实际得到的磁化强度为
$M=n\langle\mu_z\rangle$，磁化强度与饱和磁化强度之比是有用的量：
\begin{equation*}
\frac{M}{M_{\mathrm{s}}} = \frac{\langle\mu_z\rangle}{\mu} \approx \frac{y}{3} = \frac{\mu B}{3k_{\mathrm{B}}T}\tag{2.23}
\end{equation*}
再利用小场下成立的 $\chi=M/H\approx\mu_0M/B$，\fn{9}{小场时 $\chi\ll1$，故 $B\approx\mu_0H$。}得
\begin{equation*}
\chi = \frac{n\mu_0\mu^{2}}{3k_{\mathrm{B}}T}.\tag{2.24}
\end{equation*}
这表明磁化率与温度成反比，即居里定律（Curie's law，以其发现者 Pierre Curie 命名）。
\fn{10}{人们常写 $\chi=C_{\text{Curie}}/T$，$C_{\text{Curie}}$ 为居里常数。}

\subsection{$J=\frac{1}{2}$ 的顺磁性}

现在用量子系统重做上述计算：经典磁矩换成 $J=\frac{1}{2}$ 的量子自旋。磁矩 $z$
分量只有两个可能值 $m_J=\pm\frac{1}{2}$：或平行、或反平行于 $\mathbf{B}$。磁矩为
$\mp\mu_{\mathrm{B}}$（取 $g=2$），相应能量 $\pm\mu_{\mathrm{B}}B$（两个解示于
图 2.6）。于是
\begin{align*}
\langle g\mu_{\mathrm{B}}m_J\rangle &= \frac{-\mu_{\mathrm{B}}e^{\mu_{\mathrm{B}}B/k_{\mathrm{B}}T}+\mu_{\mathrm{B}}e^{-\mu_{\mathrm{B}}B/k_{\mathrm{B}}T}}{e^{\mu_{\mathrm{B}}B/k_{\mathrm{B}}T}+e^{-\mu_{\mathrm{B}}B/k_{\mathrm{B}}T}}\tag{2.25}\\
&= \mu_{\mathrm{B}}\tanh\left(\frac{\mu_{\mathrm{B}}B}{k_{\mathrm{B}}T}\right)\tag{2.26}
\end{align*}
写 $y=\mu_{\mathrm{B}}B/k_{\mathrm{B}}T=g\mu_{\mathrm{B}}JB/k_{\mathrm{B}}T$（此处
$J=\frac{1}{2}$，$g=2$），便有
\begin{equation*}
\frac{M}{M_{\mathrm{s}}} = \frac{\langle m_J\rangle}{J} = \tanh y.\tag{2.27}
\end{equation*}
这一函数不同于朗之万函数，但形状相当接近（见图 2.7）。小场下
$\tanh(\mu_{\mathrm{B}}B/k_{\mathrm{B}}T)\approx\mu_{\mathrm{B}}B/k_{\mathrm{B}}T$，
且
\begin{equation*}
\chi = \frac{n\mu_0\mu_{\mathrm{B}}^{2}}{k_{\mathrm{B}}T}.\tag{2.28}
\end{equation*}

式 2.27 也可以用另一种方法非常高效地导出。配分函数 $Z$ 是按简并度加权的玻尔兹曼
概率之和。单个自旋的配分函数为
\begin{equation*}
Z = \mathrm{e}^{\mu_{\mathrm{B}}B/k_{\mathrm{B}}T} + \mathrm{e}^{-\mu_{\mathrm{B}}B/k_{\mathrm{B}}T} = 2\cosh\left(\frac{\mu_{\mathrm{B}}B}{k_{\mathrm{B}}T}\right),\tag{2.29}
\end{equation*}
\mnote{关于 Z、F 以及 $M=-(\partial F/\partial B)_T$ 这类表达式，详见附录 E。}
由 $F=-k_{\mathrm{B}}T\ln Z$ 可得亥姆霍兹自由能；对单位体积 $n$ 个自旋，
\begin{equation*}
F = nk_{\mathrm{B}}T\ln\left[2\cosh\left(\frac{\mu_{\mathrm{B}}B}{k_{\mathrm{B}}T}\right)\right].\tag{2.30}
\end{equation*}
磁化强度由 $M=-(\partial F/\partial B)_T$ 给出，同样得到
\begin{equation*}
\frac{M}{M_{\mathrm{s}}} = \tanh\left(\frac{\mu_{\mathrm{B}}B}{k_{\mathrm{B}}T}\right),\tag{2.31}
\end{equation*}
与式 2.27 一致。

\begin{figure}
\centering
\includegraphics[width=0.42\textwidth]{fig2_06.pdf}
\caption{自旋 $\frac{1}{2}$ 磁矩的能量随磁场的变化。}
\end{figure}

\begin{figure}
\centering
\includegraphics[width=0.5\textwidth]{fig2_07.pdf}
\caption{自旋 $\frac{1}{2}$ 顺磁体的磁化强度遵循 $\tanh y$；$y$ 小时
$\tanh y\approx y$，如图中在原点附近与曲线相切的直线所示。}
\end{figure}

这一方法还可用于导出该模型的其他热力学量（见习题 2.4），结果按 $k_{\mathrm{B}}T/\mu_{\mathrm{B}}B$
绘于图 2.8。图 2.8(a) 与图 2.7 信息相同，只是横轴反向——因为要理解材料的热性质，
我们真正关心的是固定磁场下升温的效果。样品升温时磁矩随机化、磁化强度下降，但同时
能量密度 $E=-M_{\mathrm{s}}B$ 上升（见图 2.8(b)）。$T\to\infty$ 时能量为零：此时
磁矩相对外场完全随机，能量得失相抵。降温对应能量降低（这一点在 2.6 节还会提到）。

热容 $C=(\partial E/\partial T)_B$ 在 $k_{\mathrm{B}}T\sim\mu_{\mathrm{B}}B$ 附近
有一个宽峰，称为肖特基反常（Schottky anomaly，见图 2.8(c)）。其成因是：在这个
温度附近，热运动恰好能够激发系统两个态之间的跃迁。极低温下系统难以改变能量——
没有足够能量激发离开基态的跃迁，所有自旋被“冻住”，都与磁场排齐；极高温下系统也
难以改变能量——两个态被同等占据。两者之间出现极大。因此热容峰可以有效地提示
“有有趣的事情正在发生”。但要注意：肖特基反常并不是相变对应的那种尖峰，而是一个
平滑、宽缓的极大。

熵 $S=-(\partial F/\partial T)_B$ 随温度上升而增大（见图 2.8(d)），符合预期——
它反映自旋的无序程度。反之，降温对应有序化、熵减少。这一事实对磁制冷技术非常有用，
将在 2.6 节描述。

下一节考虑总角动量量子数为任意 $J$ 的一般顺磁体：前面经典与量子两种情形都作为
特例包含在内。

\begin{figure}
\centering
\includegraphics[width=0.42\textwidth]{fig2_08.pdf}
\caption{单位体积含 $n$ 个无相互作用自旋 $\frac{1}{2}$ 离子的顺磁盐的（a）磁化强度
$M$（以饱和磁化强度归一）、（b）能量 $E$、（c）热容 $C$（恒定外磁场下）与
（d）熵 $S$ 随 $k_{\mathrm{B}}T/\mu_{\mathrm{B}}B$ 的变化。$E$、$C$、$S$ 按单位体积
顺磁盐给出。}
\end{figure}

\subsection{布里渊函数}

现在推导 $J$ 取任意整数或半整数的一般情形。前面各情形（$J=\frac{1}{2}$ 与
$J=\infty$）的许多特征在这里依然成立，例如磁场使磁矩趋于排齐、温度使它们趋于
无序。配分函数为
\begin{equation*}
Z = \sum_{m_J=-J}^{J}\exp(m_J g_J\mu_{\mathrm{B}}B/k_{\mathrm{B}}T).\tag{2.32}
\end{equation*}
写 $x=g_J\mu_{\mathrm{B}}B/k_{\mathrm{B}}T$，有
\begin{equation*}
\langle m_J\rangle = \frac{\sum_{m_J=-J}^{J}m_J\,\mathrm{e}^{m_Jx}}{\sum_{m_J=-J}^{J}\mathrm{e}^{m_Jx}} = \frac{1}{Z}\frac{\partial Z}{\partial x},\tag{2.33}
\end{equation*}
于是
\begin{equation*}
M = ng_J\mu_{\mathrm{B}}\langle m_J\rangle = \frac{ng_J\mu_{\mathrm{B}}}{Z}\frac{\partial Z}{\partial B}\frac{\partial B}{\partial x} = nk_{\mathrm{B}}T\frac{\partial\ln Z}{\partial B}.\tag{2.34}
\end{equation*}
配分函数 $Z$ 是首项 $a=\mathrm{e}^{-Jx}$、公比 $r=\mathrm{e}^{x}$ 的几何级数，
可用熟知的公式求和：
\begin{equation*}
a+ar+ar^{2}+\cdots+ar^{M-1} = \sum_{j=1}^{M}ar^{j-1} = \frac{a(1-r^{M})}{1-r},\tag{2.35}
\end{equation*}
$M$ 是级数项数，此处 $M=2J+1$。稍加整理得
\begin{equation*}
Z = \frac{\sinh[(2J+1)\frac{x}{2}]}{\sinh[\frac{x}{2}]}\tag{2.36}
\end{equation*}
作代换
\begin{equation*}
y = xJ = g_J\mu_{\mathrm{B}}JB/k_{\mathrm{B}}T,\tag{2.37}
\end{equation*}
得
\begin{equation*}
M = M_{\mathrm{s}}\,\mathrm{B}_J(y)\tag{2.38}
\end{equation*}
其中饱和磁化强度
\begin{equation*}
M_{\mathrm{s}} = ng_J\mu_{\mathrm{B}}J\tag{2.39}
\end{equation*}
而 $\mathrm{B}_J(y)$ 是布里渊函数（Brillouin function）：
\bio{Léon Brillouin\\（1889--1969）}
\begin{equation*}
\mathrm{B}_J(y) = \frac{2J+1}{2J}\coth\left(\frac{2J+1}{2J}y\right) - \frac{1}{2J}\coth\frac{y}{2J}.\tag{2.40}
\end{equation*}
该函数对不同 $J$ 绘于图 2.9。布里渊函数具有恰当的极限：$J\to\infty$ 时退化为朗之万
函数
\begin{equation*}
\mathrm{B}_{\infty}(y) = L(y),\tag{2.41}
\end{equation*}
$J=\frac{1}{2}$ 时退化为 $\tanh$ 函数
\begin{equation*}
\mathrm{B}_{1/2}(y) = \tanh(y).\tag{2.42}
\end{equation*}
即回到前几节讨论过的两种情形。

$y$ 的典型量级可估计如下：取 $J=\frac{1}{2}$、$g_J=2$、$B=1$~T，室温下
$y\sim2\times10^{-3}$。因此除极低温和/或极强磁场外，实验情形都对应 $y\ll1$（因而
$\chi\ll1$）。$y$ 小时利用 $\coth y$ 的麦克劳林展开可得
\begin{equation*}
\mathrm{B}_J(y) = \frac{(J+1)y}{3J} + O(y^{3}).\tag{2.43}
\end{equation*}
低场下磁化率为
\begin{equation*}
\chi = \frac{M}{H}\approx\frac{\mu_0M}{B} = \frac{n\mu_0\mu_{\mathrm{eff}}^{2}}{3k_{\mathrm{B}}T}\tag{2.44}
\end{equation*}
形如经典的居里定律。\fn{11}{即 $\chi=C_{\text{Curie}}/T$，其中 $C_{\text{Curie}}=n\mu_0g_J^{2}J(J+1)/3k_{\mathrm{B}}$ 为居里常数。}于是测量 $\chi$ 便可推定有效磁矩
$\mu_{\mathrm{eff}}$：
\begin{equation*}
\mu_{\mathrm{eff}} = g_J\mu_{\mathrm{B}}\sqrt{J(J+1)}\tag{2.45}
\end{equation*}
其中
\begin{equation*}
g_J = \frac{3}{2} + \frac{S(S+1)-L(L+1)}{2J(J+1)}\tag{2.46}
\end{equation*}
常数 $g_J$ 称为 Landé $g$ 值（见附录 C）。
\bio{Alfred Landé\\（1888--1975）}

\begin{figure}
\centering
\includegraphics[width=0.55\textwidth]{fig2_09.pdf}
\caption{磁矩量子数为 $J$ 的顺磁体的磁化强度遵循布里渊函数 $\mathrm{B}_J(y)$，
图中对不同 $J$ 绘出：$J=\frac{1}{2},1,\frac{3}{2},2,\frac{5}{2},\ldots$ 及 $J=\infty$。}
\end{figure}

居里定律的 $\chi\propto1/T$ 意味着 $1/\chi$ 对 $T$ 作图是直线、$\chi T$ 对 $T$ 作图
为常数（见图 2.10）。后续章节考虑相互作用的作用时，记住这几点会有用。还须注意：
磁化率是在外磁场趋于零的极限下取值，由式 2.44 以
$\mu_{\mathrm{eff}}=g_J\mu_{\mathrm{B}}\sqrt{J(J+1)}$ 给出；而在强场下磁化强度饱和
为 $M_{\mathrm{s}}$，由式 2.39 相当于每离子磁矩 $g_J\mu_{\mathrm{B}}J$。这两个数值
$\mu_{\mathrm{eff}}=g_J\mu_{\mathrm{B}}\sqrt{J(J+1)}$ 与 $M_{\mathrm{s}}/n=g_J\mu_{\mathrm{B}}J$
一般不等，仅当 $J\to\infty$（经典极限）时才相等。

\begin{figure}
\centering
\includegraphics[width=0.32\textwidth]{fig2_10.pdf}
\caption{居里定律 $\chi\propto1/T$（a）；$1/\chi$ 对 $T$ 作图为直线（b）；$\chi T$
对 $T$ 作图为常数（c）。}
\end{figure}

\subsection{范弗莱克顺磁性}

若基态 $|0\rangle$ 的 $J=0$，则没有顺磁效应，因为
\begin{equation*}
\langle 0|\hat{\boldsymbol{\mu}}|0\rangle = g_J\mu_{\mathrm{B}}\langle 0|\hat{\mathbf{J}}|0\rangle = 0.\tag{2.47}
\end{equation*}
这意味着加磁场不改变系统基态能量，因而没有顺磁磁化率。但这一结论只在一阶微扰论
中正确。二阶微扰论预言基态能量 $E_0$ 会变——因为它计入了 $J\neq0$ 的激发态混入。
$J=0$ 离子的基态能量变化为
\begin{equation*}
\Delta E_0 = \sum_n\frac{|\langle 0|(\mathbf{L}+g\mathbf{S})\cdot\mathbf{B}|n\rangle|^{2}}{E_0-E_n} + \frac{e^{2}}{8m_{\mathrm{e}}}\sum_i(\mathbf{B}\times\mathbf{r}_i)^{2},\tag{2.48}
\end{equation*}
第二项来自抗磁性，第一项求和遍及系统所有激发态。磁化率于是为
\begin{equation*}
\chi = \frac{N}{V}\left(2\mu_{\mathrm{B}}^{2}\sum_n\frac{|\langle 0|(L_z+gS_z)|n\rangle|^{2}}{E_n-E_0} - \frac{e^{2}\mu_0}{6m_{\mathrm{e}}}\sum_{i=1}^{Z}\langle r_i^{2}\rangle\right),\tag{2.49}
\end{equation*}
第一项为正（因 $E_n>E_0$），称为范弗莱克顺磁性（Van Vleck paramagnetism）；
第二项为负，即前面讨论过的常规抗磁磁化率（式 2.15）。与抗磁性一样，范弗莱克
顺磁性既小、又不依赖温度。
\bio{John H. van Vleck\\（1899--1980）}

\section{离子的基态与洪特规则}

典型的原子含有很多电子，其中大部分位于没有净角动量的填满壳层中。但可能存在未满
壳层，其中电子可以组合出非零的自旋与轨道角动量。1.3.4 节已经看到两个自旋
$\frac{1}{2}$ 电子如何组合成自旋为 0 或 1 的联合体。在原子中，未满壳层所有电子的
自旋角动量可以全部组合在一起，轨道角动量亦然。于是原子具有总轨道角动量
$\hbar\mathbf{L}$ 与总自旋角动量 $\hbar\mathbf{S}$。轨道与自旋角动量可以按
\begin{equation*}
(2L+1)(2S+1)\tag{2.50}
\end{equation*}
种方式组合：这是 $\mathbf{L}$ 的 $z$ 分量的取值数（级数 $-L,-L+1,\ldots,L-1,L$
共 $(2L+1)$ 项）乘以 $\mathbf{S}$ 的 $z$ 分量的取值数（同理为 $(2S+1)$）。这些以
不同方式组合未满壳层电子角动量（自旋与轨道）的组态具有不同的能量。原因在于：自旋
角动量的选择影响波函数的空间部分，轨道角动量影响电子绕核的运动方式——两者都影响
电子相互回避的好坏，从而影响静电排斥能。下面说明如何找到能量最低的组态。

\subsection{精细结构}

迄今我们一直把自旋与轨道角动量分开处理，因为它们彼此独立。然而它们确实通过
自旋--轨道相互作用（spin--orbit interaction，见附录 C）弱耦合，它对具有确定 L、S
的态是一个微扰。因此 L 与 S 不再分别守恒，但总角动量 $\mathbf{J}=\mathbf{L}+\mathbf{S}$
守恒。若把相对论效应作为微扰处理（通常可以），则 $\mathbf{L}^{2}=L(L+1)$ 与
$\mathbf{S}^{2}=S(S+1)$ 可视为守恒。于是具有给定 L、S 的态分裂成若干不同 J 的能级，
称为精细结构（fine structure）。J 取 $|L-S|$ 到 $L+S$ 之间的值。由 J 的定义
（式 2.17），
\begin{equation*}
\mathbf{J}^{2} = \mathbf{L}^{2}+\mathbf{S}^{2}+2\mathbf{L}\cdot\mathbf{S},\tag{2.51}
\end{equation*}
自旋--轨道相互作用取 $\lambda\mathbf{L}\cdot\mathbf{S}$ 的形式（见附录 C），$\lambda$
为常数，其能量期望值为
\begin{equation*}
\langle\lambda\mathbf{L}\cdot\mathbf{S}\rangle = \frac{\lambda}{2}[J(J+1)-L(L+1)-S(S+1)].\tag{2.52}
\end{equation*}
原子的能量主要由静电考虑通过 S、L 的取值决定，因此能量本征态可以用 S 和 L 标记。
在无自旋--轨道相互作用时，J 在 $|L-S|$ 到 $L+S$ 中的具体取值并不影响能量。每个
能级是 $(2S+1)(2L+1)$ 个态组成的多重态（multiplet）。把自旋--轨道相互作用作为
微扰加入后，多重态分裂成以 J 标记的不同精细结构能级；每个能级自身又有 $2J+1$ 的
简并，故加磁场后可按不同 $m_J$ 进一步分裂。精细结构能级的分裂遵循朗德间隔定则
（Land\'e interval rule）：给定多重态中相邻能级 $E(J)$ 与 $E(J-1)$ 的间隔为
\begin{align*}
E(J)-E(J-1) &= \frac{\lambda}{2}[J(J+1)-L(L+1)-S(S+1)]\\
&\quad -\frac{\lambda}{2}[(J-1)J-L(L+1)-S(S+1)]\\
&= \lambda J.\tag{2.53}
\end{align*}
即（考虑 $J-1$ 与 $J$ 两能级的间隔）分裂正比于 $J$。

\begin{figure}
\centering
\includegraphics[width=0.5\textwidth]{fig2_11.pdf}
\caption{角动量组合 $L=3$、$S=\frac{3}{2}$ 给出从 $\frac{3}{2}$ 到 $\frac{9}{2}$ 的
各个 $J$ 值。}
\end{figure}

\begin{example}{2.1}
图 2.11 以 $L=3$、$S=\frac{3}{2}$ 为例演示这种角动量组合。此时
$(2S+1)(2L+1)=28$；由图可见，若考虑 J 从 $|L-S|=\frac{3}{2}$ 到 $L+S=\frac{9}{2}$
的各个态，由于每个 J 态简并度为 $2J+1$，总共也是 28 个态。这形象地说明了
\begin{equation*}
\sum_{J=|L-S|}^{L+S}(2J+1) = (2L+1)(2S+1).\tag{2.54}
\end{equation*}
\end{example}

角动量量子数的组合显然非常多；对一个特定离子，哪一个是基态？

\subsection{洪特规则}

使能量取极小的角动量量子数组合可以用洪特规则（Hund's rules）估计。\fn{12}{洪特规则只适用于基态组态，对最低能级之上各能级的次序不作任何断言；它还假定只有一个未满支壳层。}\bio{Friederich Hund\\（1896--1997）}这三条经验规则按重要性递减排列：先满足第一条，在此前提下再尽量满足第二条，依此类推。

（1）安排电子波函数使 S 最大。由于泡利不相容原理阻止自旋平行的电子占据同一位置，
库仑能得以减小，从而降低电子间的静电排斥。

（2）在第一条已定的波函数基础上使 L 最大。这同样降低能量：同向旋转轨道上的电子
能更有效地相互回避，从而减小库仑排斥。

（3）最后用 $J=|L-S|$（壳层填充不足半满）或 $J=|L+S|$（超过半满）定出 J。第三条
规则源于使自旋--轨道能量取极小的尝试。注意它只在一定条件下适用：如下一章所示，
在许多体系（过渡金属离子是很好的例子）中，自旋--轨道能量不如晶体场等其他能量项
重要，此时洪特第三条规则被违反；而对稀土离子，如下所示，第三条规则非常好用。

求出 S、L、J 之后，基态可用形如 ${}^{2S+1}L_J$ 的谱项符号（term symbol）概括。
其中 L 不写成数字，而按下列序列用字母表示：
\begin{center}
\begin{tabular}{lcccccccc}
\toprule
$L$ & 0 & 1 & 2 & 3 & 4 & 5 & 6 & $\cdots$\\
\midrule
 & S & P & D & F & G & H & I & $\cdots$\\
\bottomrule
\end{tabular}
\end{center}
$2S+1$ 是自旋多重度（spin multiplicity）。

\begin{example}{2.2}
以稀土离子 $\mathrm{Dy^{3+}}$（外壳层 $4f^{9}$）为例：f 电子 $l=3$，为满足洪特
第一条规则，$2l+1=7$ 个自旋向上，剩下 2 个自旋向下（见图 2.12）。于是
$S=7\times\frac{1}{2}-2\times\frac{1}{2}=\frac{5}{2}$（自旋简并 $2S+1=6$）。自旋
向上的电子对轨道角动量没有净贡献，因此只需使其余 2 个自旋向下的电子的轨道贡献最大：
$L=3+2=5$，对应字母 H。壳层超过半满，故 $J=|5+\frac{5}{2}|=\frac{15}{2}$。谱项
符号为 ${}^{6}\mathrm{H}_{15/2}$。
\end{example}

\begin{figure}
\centering
\includegraphics[width=0.35\textwidth]{fig2_12.pdf}
\caption{$\mathrm{Dy^{3+}}$ 的基态。}
\end{figure}

洪特规则预言基态，但对激发态及其与基态的接近程度不置一词。因此它使我们能在
“只有基态被占据”的假定下估计离子的磁矩。

本章和下一章都要把这些预言与实验比较，我们将集中于含 3d 与 4f 离子的化合物——
它们在许多磁性体系中都很重要。3d 元素是第一行过渡金属（Sc--Zn），4f 元素即镧系
元素或稀土元素（La--Lu），见图 2.13 的周期表。

\begin{figure}
\centering
\includegraphics[width=0.85\textwidth]{fig2_13.pdf}
\caption{周期表。3d 元素（Sc--Zn）与 4f 元素（La--Lu）以阴影标出。每个元素旁的数字
是原子序数 Z（核内质子数）。}
\end{figure}

洪特规则用于 3d 与 4f 离子的结果示于图 2.14（注意：如下一章详述，4f 离子的洪特
规则预言比 3d 离子更符合实验）。\fn{13}{本章只讨论自由原子或离子。把原子放入晶体环境中情况就会改变；对 3d 离子变化相当大，见第三章。}在每个组中段 S 升至极大；L 与
J 的极大大致出现在四分之一与四分之三填充处——J 的两个极大不对称，反映了壳层不足
半满与超过半满时规则的不同。

\begin{figure}
\centering
\includegraphics[width=0.48\textwidth]{fig2_14.pdf}
\caption{按洪特规则得到的 3d 与 4f 离子的 S、L、J。图中 $n$ 是支壳层（3d 或 4f）
中的电子数。}
\end{figure}

由式 2.44，测量磁化率可以推定有效磁矩。以玻尔磁子为单位：
\begin{equation*}
p = \mu_{\mathrm{eff}}/\mu_{\mathrm{B}},\tag{2.55}
\end{equation*}
洪特规则的预言（用式 2.45）是 $p=g_J[J(J+1)]^{\frac{1}{2}}$。如表 2.2 所示，对
固体中的 4f 离子，这一预言与实测 $p=\mu_{\mathrm{eff}}/\mu_{\mathrm{B}}$ 通常符合
极好。Sm 与 Eu 出现偏差，原因是它们存在 J 不同于基态的低激发态，这些态
因距基态很近而也有显著布居，使 $p$ 偏离基态值。4f 离子总体符合良好，但下一章会
看到 3d 离子的符合程度差得多，这是局域晶体环境影响所致。这一效应对 4f 离子不
重要，因为部分填充的 4f 壳层深藏于离子内部、处于填满的 5s 与 5p 壳层之下。

\begin{table}
\centering
\caption{用洪特规则得到的 4f 离子磁基态。对每个离子列出壳层组态及预言的基态
S、L、J；并给出由洪特规则计算的 $p=\mu_{\mathrm{eff}}/\mu_{\mathrm{B}}=g_J[J(J+1)]^{1/2}$。
下一列是实验值 $p_{\mathrm{exp}}$：除 Sm、Eu 外符合极好。实验值来自顺磁盐在
$k_{\mathrm{B}}T\gg E_{\mathrm{CEF}}$ 温度下的磁化率测量，$E_{\mathrm{CEF}}$ 为
晶体场能量。}
\small
\begin{tabular}{lccccccc}
\toprule
离子 & 壳层 & $S$ & $L$ & $J$ & 谱项 & $p$ & $p_{\mathrm{exp}}$\\
\midrule
$\mathrm{Ce^{3+}}$ & $4f^{1}$ & $\frac{1}{2}$ & 3 & $\frac{5}{2}$ & ${}^{2}F_{5/2}$ & 2.54 & 2.51\\
$\mathrm{Pr^{3+}}$ & $4f^{2}$ & 1 & 5 & 4 & ${}^{3}H_{4}$ & 3.58 & 3.56\\
$\mathrm{Nd^{3+}}$ & $4f^{3}$ & $\frac{3}{2}$ & 6 & $\frac{9}{2}$ & ${}^{4}I_{9/2}$ & 3.62 & 3.3--3.7\\
$\mathrm{Pm^{3+}}$ & $4f^{4}$ & 2 & 6 & 4 & ${}^{5}I_{4}$ & 2.68 & --\\
$\mathrm{Sm^{3+}}$ & $4f^{5}$ & $\frac{5}{2}$ & 5 & $\frac{5}{2}$ & ${}^{6}I_{5/2}$ & 0.85 & 1.74\\
$\mathrm{Eu^{3+}}$ & $4f^{6}$ & 3 & 3 & 0 & ${}^{7}F_{0}$ & 0.0 & 3.4\\
$\mathrm{Gd^{3+}}$ & $4f^{7}$ & $\frac{7}{2}$ & 0 & $\frac{7}{2}$ & ${}^{8}S_{7/2}$ & 7.94 & 7.98\\
$\mathrm{Tb^{3+}}$ & $4f^{8}$ & 3 & 3 & 6 & ${}^{7}F_{6}$ & 9.72 & 9.77\\
$\mathrm{Dy^{3+}}$ & $4f^{9}$ & $\frac{5}{2}$ & 5 & $\frac{15}{2}$ & ${}^{6}H_{15/2}$ & 10.63 & 10.63\\
$\mathrm{Ho^{3+}}$ & $4f^{10}$ & 2 & 6 & 8 & ${}^{5}I_{8}$ & 10.60 & 10.4\\
$\mathrm{Er^{3+}}$ & $4f^{11}$ & $\frac{3}{2}$ & 6 & $\frac{15}{2}$ & ${}^{4}I_{15/2}$ & 9.59 & 9.5\\
$\mathrm{Tm^{3+}}$ & $4f^{12}$ & 1 & 5 & 6 & ${}^{3}H_{6}$ & 7.57 & 7.61\\
$\mathrm{Yb^{3+}}$ & $4f^{13}$ & $\frac{1}{2}$ & 3 & $\frac{7}{2}$ & ${}^{2}F_{7/2}$ & 4.53 & 4.5\\
$\mathrm{Lu^{3+}}$ & $4f^{14}$ & 0 & 0 & 0 & ${}^{1}S_{0}$ & 0 & 0\\
\bottomrule
\end{tabular}
\end{table}

本章我们把每个原子的磁矩与局域在该原子上的电子相联系。金属中的传导电子并不局域
于特定原子，而是处于能带态、遍布整个材料。传导电子既表现抗磁效应也表现顺磁效应，
这将在第七章讨论。

\subsection{L--S 耦合与 j--j 耦合}

以上原子模型假定 L--S 耦合（L--S coupling，又称 Russell--Saunders 耦合）：自旋--
轨道相互作用是弱微扰，主要能量项由决定 L 与 S 取值的静电相互作用给出——即先把
各电子的轨道与自旋角动量分别组合起来；直到知道整个原子的总轨道与总自旋角动量
之后，才把自旋--轨道相互作用作为弱微扰加入，把每个谱项分裂成以 J 标记的精细结构
能级。

对高原子序数 Z 的原子这不再适用，因为自旋--轨道相互作用能正比于
$Z^{4}$\fn{14}{$Z^{4}$ 律只对类氢原子成立；对中性原子没有这么简单，但接近 $Z^{2}$。}（见附录 C），在这些原子中不能当作小微扰。更好的方案称为 j--j 耦合（j--j coupling）：
此时自旋--轨道相互作用是主导能量，先把每个电子自身的自旋与轨道角动量耦合起来，
随后较弱的静电效应再耦合各电子的总角动量。

\begin{example}{2.3}
以碳（C）与铅（Pb）为例。二者外壳层都有两个 p 电子，但 Pb 在周期表中比 C 靠后得
多，其 $Z^{4}$ 大近五个数量级。碳（组态 $2p^{2}$）可用 L--S 耦合处理：由洪特规则，
S=1（规则一“S 最大”，两电子 $\frac{1}{2}+\frac{1}{2}$）、L=1（规则二“L 最大”，
p 电子一个可取 $m_l=1$，另一个只能取 $m_l=0$）、J=0（规则三，不足半满，
$J=|1-1|$），故 C 基态谱项符号为 ${}^{3}P_{0}$。

铅（组态 $6p^{2}$）中自旋--轨道相互作用占主导，更适合 j--j 耦合。每个电子
$s=\frac{1}{2}$、$l=1$，组合 j 可为 $\frac{1}{2}$ 或 $\frac{3}{2}$，取决于自旋--
轨道能的符号；实际给出 $j=\frac{1}{2}$。两个电子都有 $j=\frac{1}{2}$，经较弱的
静电效应耦合成可为 0 或 1 的总 J。由泡利不相容原理基态为 J=0（两个电子若量子数
全同就必须处于反对称态而分开），故基态的好量子数为 $j_1=\frac{1}{2}$、
$j_2=\frac{1}{2}$、J=0。注意此时 L 与 S 不是好量子数。（好量子数的定义见附录 C。）
\end{example}

\section{绝热退磁}

本节描述一种把样品冷却到极低温度的技术，它要用到我们关于顺磁材料的一些想法，
并相当清楚地带出熵的概念。考虑含 $N$ 个独立磁矩的顺磁盐样品。不加磁场时磁矩方向
随机（我们假定它们彼此无相互作用），系统没有净磁化强度。加磁场则使磁矩趋于排齐、
产生磁化强度。升温减小磁化强度，加场增大磁化强度——这体现于磁化强度是 $B/T$ 的
函数（见式 2.37 与 2.38）。

极高温下所有磁矩方向完全随机，净磁化强度为零。热能 $k_{\mathrm{B}}T$ 太大，所有
态被同等占据，无论其能量是否有利。若磁矩 $J=\frac{1}{2}$，只能平行或反平行于磁场：
排列上下磁矩的方式数 $W=2^{N}$。于是磁矩对熵 S 的贡献为
\begin{equation*}
S = k_{\mathrm{B}}\ln W = Nk_{\mathrm{B}}\ln 2.\tag{2.56}
\end{equation*}
一般 $J>\frac{1}{2}$ 时 $W=(2J+1)^{N}$，熵为
\begin{equation*}
S = Nk_{\mathrm{B}}\ln(2J+1).\tag{2.57}
\end{equation*}

计算较低温度下的熵需要换个方法。离子总角动量 $z$ 分量取 $m_J$ 的概率
$p(m_J)$ 正比于玻尔兹曼因子 $\exp(-m_Jg\mu_{\mathrm{B}}B/k_{\mathrm{B}}T)$：
\begin{equation*}
p(m_J) = \frac{1}{Z}\exp(m_J g\mu_{\mathrm{B}}B/k_{\mathrm{B}}T)\tag{2.58}
\end{equation*}
$Z$ 是配分函数（式 2.32）。系统的有序程度由熵 S 描述；既然现在概率化地处理系统，
就采用表达式
\begin{equation*}
S = -Nk_{\mathrm{B}}\sum_{m_J}p(m_J)\ln p(m_J).\tag{2.59}
\end{equation*}
熵的表达式也可以这样得到：先由 $F=-Nk_{\mathrm{B}}T\ln Z$ 算出亥姆霍兹自由能 F，
再用 $S=-(\partial F/\partial T)_B$。

考察式 2.59 的后果。无外场或高温下系统完全无序，各 $m_J$ 等概率
$p(m_J)=1/(2J+1)$，熵退化为
\begin{equation*}
S = Nk_{\mathrm{B}}\ln(2J+1)\tag{2.60}
\end{equation*}
与式 2.57 一致。温度降低时低能量态越来越可几；磁矩沿外场排列的程度（磁化强度）
增加，熵下降。极低温下所有磁矩都与磁场排齐以降低能量——此时系统只有一种排列方式
（全部自旋排齐），$W=1$，$S=0$。

磁制冷样品的原理如下。先用液氦把顺磁体冷却到一个低的起始温度，然后分两步进行
磁制冷（见图 2.15）。

\begin{figure}
\centering
\includegraphics[width=0.6\textwidth]{fig2_15.pdf}
\caption{顺磁盐的熵随温度的函数关系，对应零与某个最大值 $B_b$ 之间的若干外加磁场。
把顺磁盐从温度 $T_{\mathrm{i}}$ 磁制冷到 $T_{\mathrm{f}}$ 分两步：先等温磁化，
在恒温 $T_{\mathrm{i}}$ 下把磁场从 0 增至 $B_b$（$a\to b$）；再绝热退磁
（$b\to c$）。$S(T)$ 曲线的计算取 $J=\frac{1}{2}$（见式 2.76），并加了一项
$\propto T^{3}$ 以模拟晶格振动的熵。$B=0$ 的曲线实际对应小而非零的 B，以模拟小的
剩余场。}
\end{figure}

第一步是等温磁化。顺磁体的能量因磁矩沿场排列而降低，故给定温度下增强外场可以增加
排齐程度。这一步等温进行（图 2.15 中 $a\to b$）：让样品与液氦池保持热接触（氦在
大气压下的沸点为 4.2 K，减压还可低于 4.2 K）。样品温度不变，氦池吸收样品能量与熵
下降时释放的热。热接触通常由样品室中的低压氦气提供——它在样品与室壁之间导热，
室本身浸在氦池中。（这种气体常称“交换”气体，因为它让样品与热池交换热量。）

第二步是使样品与氦池热隔离（抽走交换气体），然后缓慢把磁场降至零——慢到过程是
准静态的、熵保持不变。这一步称为绝热退磁（adiabatic demagnetization，图 2.15 中
$b\to c$），它降低系统温度。绝热退磁期间样品的总熵不变：磁矩的熵增加（场减弱时
磁矩随机化），恰好被声子（晶格振动）的熵减少（样品降温）所平衡。熵就这样在声子
与自旋之间交换。

这种制冷方法有极限吗？初看之下，$B=0$ 时熵在一切 $T>0$ 似乎都是
$S=Nk_{\mathrm{B}}\ln(2J+1)$，只在绝对零度才降为零，那么绝热退磁似乎可以一路制冷
到绝对零度。然而真实系统中总存在磁矩间相互作用导致的微小剩余内场，使熵在略高于
绝对零度处就提前跌落（见图 2.15）。这个场的大小限制了盐能冷却到的最低温度。某些
剩余内场极小的盐可以达到几个毫开尔文。

\section{核自旋}

原子中具有磁矩的不只是电子。原子核常因核角动量而有非零自旋。每个核有一个称为核
自旋量子数的量子数 $I$，代表以 $\hbar$ 为单位的核总角动量。但核磁矩非常小，因为
其大小与所涉粒子的质量成反比：核磁矩典型地比电子磁矩小三个量级（大多介于
$10^{-3}$ 与 $10^{-4}\mu_{\mathrm{B}}$ 之间）。磁矩之小，加上相邻原子核之间不存在
强相互作用，使核自旋体系在常温实验室条件下无法有序。

\begin{table}
\centering
\caption{一些常见核自旋的性质。首先列出中子（n）、质子（p）、氘核（d）与氚核（t）。
Z 是原子序数（质子数），N 是中子数，核 X 记作 ${}^{A}X$，$A=Z+N$ 为质量数。I 是核
自旋量子数，$\mu$ 是以核磁子（$\mu_{\mathrm{N}}$）量度的核磁矩，$g_I$ 是核 $g$
因子。最后一列是这些磁矩在 1 T 磁场中的进动频率 $\nu$。}
\small
\begin{tabular}{lcccccc}
\toprule
核 & $Z$ & $N$ & $I$ & $\mu/\mu_{\mathrm{N}}$ & $g_I$ & $\nu$（B=1 T，MHz）\\
\midrule
n & 0 & 1 & $\frac{1}{2}$ & $-1.913$ & $-3.826$ & 29.17\\
p=${}^{1}$H & 1 & 0 & $\frac{1}{2}$ & 2.793 & 5.586 & 42.58\\
d=${}^{2}$H & 1 & 1 & 1 & 0.857 & 0.857 & 6.536\\
t=${}^{3}$H & 1 & 2 & $\frac{1}{2}$ & $-2.128$ & 4.255 & 32.43\\
${}^{12}$C & 6 & 6 & 0 & 0 & 0 & 0\\
${}^{13}$C & 6 & 7 & $\frac{1}{2}$ & 0.702 & 1.404 & 10.71\\
${}^{14}$N & 7 & 7 & 1 & 0.404 & 0.404 & 3.076\\
${}^{16}$O & 8 & 8 & 0 & 0 & 0 & 0\\
${}^{17}$O & 8 & 9 & $\frac{5}{2}$ & $-1.893$ & $-0.757$ & 5.772\\
${}^{19}$F & 9 & 10 & $\frac{1}{2}$ & 2.628 & 5.257 & 40.05\\
${}^{31}$P & 15 & 16 & $\frac{1}{2}$ & 1.132 & 2.263 & 17.24\\
${}^{33}$S & 16 & 17 & $\frac{3}{2}$ & 0.643 & 0.429 & 3.266\\
\bottomrule
\end{tabular}
\end{table}

核磁性的单位是核磁子（nuclear magneton）$\mu_{\mathrm{N}}$：
\begin{equation*}
\mu_{\mathrm{N}} = \frac{e\hbar}{2m_{\mathrm{p}}} = 5.0508\times10^{-27}\ \mathrm{A\,m^{2}},\tag{2.61}
\end{equation*}
$m_p$ 是质子质量。（这远小于 1s 电子的磁矩——玻尔磁子
$\mu_{\mathrm{B}}=e\hbar/2m_{\mathrm{e}}=9.27\times10^{-24}$ A\,m$^2$，$m_{\mathrm{e}}$
为电子质量，见式 1.15。）核自旋量子数 I 取下列数值之一：0、$\frac{1}{2}$、1、
$\frac{3}{2}$、2、……角动量沿 $z$ 方向的分量为 $m_I$（以 $\hbar$ 为单位），$m_I$
可取
\begin{equation*}
-I, -I+1, \dots, I-1, I.\tag{2.62}
\end{equation*}
把核磁矩投影到 $z$ 方向、取 $m_I$ 最大的自旋态，其值为
\begin{equation*}
\mu = g_I\mu_{\mathrm{N}}I\tag{2.63}
\end{equation*}
$g_I$ 是核 $g$ 因子，为 1 量级的数，反映核的细致结构。若干核的 I、$g_I$ 与 $\mu$
列于表 2.3。中子与质子的磁矩都不是 $\mu_{\mathrm{N}}$ 的简单倍数——这表明每个粒子
内部结构复杂：中子与质子都由三个夸克组成。质子 $I=\frac{1}{2}$、$g_I=5.586$，磁矩
$\mu_p=1.410\times10^{-26}$ A\,m$^2$。

通常我们忽略核磁性：水中质子的磁矩并不能让水吸附在磁极上（事实上由于水分子中
电子云，水主要是弱抗磁体）。不过，借助非常灵敏的共振技术实验可以探测核磁矩。这一
实验方法称为核磁共振（nuclear magnetic resonance），将在 3.2.1 节讨论。

\begin{example}{2.4}
许多常见核完全没有磁矩，例如最常见的碳与氧——${}^{12}\mathrm{C}$ 与
${}^{16}\mathrm{O}$。这是因为这两种核的质子数与中子数都是偶数；核中全同粒子对
倾向于配成单态。有未成对核子的核才有非零磁矩：${}^{13}\mathrm{C}$ 有一个未成对
中子，$I=\frac{1}{2}$；${}^{14}\mathrm{N}$ 有一个未成对中子和一个未成对质子，
组合给出 I=1。虽然碳与氧最丰的同位素没有核磁矩，其次丰的同位素却有（天然碳约
1.1\% 是 ${}^{13}\mathrm{C}$，天然氧约 0.04\% 是 ${}^{17}\mathrm{O}$；二者都有
非零磁矩，见表 2.3）。
\end{example}

\section{超精细结构}

原子内部，核磁矩可以与电子磁矩发生磁相互作用，但极其微弱。由此产生的能级劈裂
比 2.5.1 节的精细结构还要小，故称超精细结构（hyperfine structure）。核与电子之间
的弱磁相互作用称为超精细相互作用（hyperfine interactions）。其起源相当复杂，
但基本原理可以这样理解：考虑核磁矩 $\boldsymbol{\mu}$ 处于磁场
$\mathbf{B}_{\text{electrons}}$ 中（由所有电子的运动与自旋产生），产生能量项
$-\boldsymbol{\mu}\cdot\mathbf{B}_{\text{electrons}}$。$\mathbf{B}_{\text{electrons}}$
预期正比于所有电子的总角动量 $\mathbf{J}$，故超精细相互作用的哈密顿量可写为
\begin{equation*}
\hat{\mathcal{H}}_{\mathrm{hf}} = A\,\mathbf{I}\cdot\mathbf{J}.\tag{2.64}
\end{equation*}
I 是核角动量，A 是可由实验确定的参数，量度超精细相互作用的强度。

对一般原子，超精细相互作用的严格形式很复杂，但对与点核相互作用的单电子原子可以
详细算出（见习题 2.8）。这里只作概述，以阐明相互作用的基本机制。

重要的磁相互作用有两类。第一类是磁偶极相互作用（4.1 节详述）。核磁矩
$\boldsymbol{\mu}_I=g_I\mu_I\mathbf{I}$ 与电子磁矩
$\boldsymbol{\mu}_{\mathrm{e}}=g_e\mu_{\mathrm{B}}\mathbf{S}$ 之间的偶极相互作用为
$\hat{\mathcal{H}}_{\text{dipole}}$：
\begin{align*}
\hat{\mathcal{H}}_{\text{dipole}} &= \frac{\mu_0}{4\pi r^{3}}\left[\boldsymbol{\mu}_{\mathrm{e}}\cdot\boldsymbol{\mu}_I - \frac{3}{r^{2}}(\boldsymbol{\mu}_{\mathrm{e}}\cdot\mathbf{r})(\boldsymbol{\mu}_I\cdot\mathbf{r})\right]\tag{2.65}\\
&= \frac{\mu_0 g_e g_I\mu_{\mathrm{B}}\mu_I}{4\pi r^{3}}\left[\mathbf{I}\cdot\mathbf{S} - \frac{3}{r^{2}}(\mathbf{S}\cdot\mathbf{r})(\mathbf{I}\cdot\mathbf{r})\right].\tag{2.66}
\end{align*}
若两个磁矩都沿 $z$ 轴排列，则 $\hat{\mathcal{H}}_{\text{dipole}}\propto(1-3\cos^{2}\theta)I_zS_z/r^{3}$，
相互作用能正比于
\begin{equation*}
\int\mathrm{d}^{3}\mathbf{r}\,|\psi(\mathbf{r})|^{2}(1-3\cos^{2}\theta)m_sm_I/r^{3},\tag{2.67}
\end{equation*}
对 s 轨道该积分为零——本质原因是偶极场在球面上平均为零，可借图 2.16 理解。

球对称的 s 轨道波函数把它平均为零。但对 $l>0$ 轨道（如 p 轨道）中的未成对电子，
积分不为零。上面的处理也只考虑了电子磁矩的自旋部分；电子的轨道磁矩同样在核处产生
磁场，给出正比于 $\mathbf{I}\cdot\mathbf{L}$ 的项。（所有这些项在习题 2.8 的处理中
都会自然出现。）

\begin{figure}
\centering
\includegraphics[width=0.5\textwidth]{fig2_16.pdf}
\caption{磁性核（设为球形，图中以阴影表示）内外磁场分布的示意图。磁场分布由绕核赤道
流动的电流（图中阴影）产生。球外的场看似来自球心处的点偶极子，对球对称体积求平均
为零——场向上与向下的机会均等。核球内部的场平均向上、对球对称体积平均不为零。
这正是费米接触相互作用的起源。}
\end{figure}

第二类机制是费米接触相互作用（Fermi contact interaction）——与上一机制相反，
它对所有非 s 轨道为零，只对 s 轨道不为零。这里考虑式 2.66 在 $r\to0$ 时的情形。
核当然不是点偶极子，而有有限体积，记为 $V$。我们进一步假设\fn{15}{这一假设只是为了让计算简单。它实际上给出正确答案，但更好的处理见习题 2.8。}核是磁化强度均匀的球
$\mathbf{M}=\boldsymbol{\mu}_I/V$。球内的磁通密度为 $\mu_0\mathbf{M}$，但须先扣除
退磁场——对球为 $\mu_0\mathbf{M}/3$（见附录 D）——剩下 $2\mu_0\mathbf{M}/3$。
于是闯入核球内的电子将感受到磁通密度
\begin{equation*}
\mathbf{B} = \frac{2\mu_0\boldsymbol{\mu}_I}{3V}.\tag{2.68}
\end{equation*}
能量代价由它乘以 $\mu_{\mathrm{e}}|\psi(0)|^{2}V$——处于核内的电子磁矩量（注意
$|\psi(0)|^{2}V$ 是在核内找到电子的概率）——得到。故接触能 $E_{\text{contact}}$ 为
\begin{equation*}
E_{\text{contact}} = \left(\frac{2\mu_0}{3}\right)2\mu_{\mathrm{B}}g_I\mu_{\mathrm{N}}\,\mathbf{I}\cdot\mathbf{S},\tag{2.69}
\end{equation*}
其中用了 $\boldsymbol{\mu}_{\mathrm{e}}=-2\mu_{\mathrm{B}}\mathbf{S}$ 与
$\boldsymbol{\mu}_I=g_I\mu_I\mathbf{I}$。电子在核内感受到的净场也可借图 2.16 理解。
只有 s 电子在 $r=0$ 处有振幅，因而只有它们表现出这一效应。费米接触相互作用引起的
s 电子超精细劈裂，通常远大于 $l>0$ 电子的磁偶极超精细劈裂。

如习题 2.8 所示，同时包含偶极与费米接触项的超精细相互作用形式相当复杂，但其要点
是：对 s 电子（$l=0$）为常数倍的 $\mathbf{I}\cdot\mathbf{S}$，对 $l>0$ 电子为
$\mathbf{I}\cdot\mathbf{N}$，其中 $\mathbf{N}=\mathbf{L}-\mathbf{S}+3(\mathbf{S}\cdot\mathbf{r})\mathbf{r}/r^{2}$。
对 s 电子 $\mathbf{J}=\mathbf{S}$，相互作用与式 2.64 一致。可以证明 $\mathbf{N}$
能投影到 $\mathbf{J}$ 上，使 $\mathbf{I}\cdot\mathbf{N}$ 的期望值正比于
$\mathbf{I}\cdot\mathbf{J}$（对单电子原子比例常数为 $L(L+1)/J(J+1)$），故式 2.64
对 $l>0$ 的电子也适用。

\begin{figure}
\centering
\includegraphics[width=0.42\textwidth]{fig2_17.pdf}
\caption{（a）$I=\frac{1}{2}$、$J=\frac{1}{2}$ 与（b）$I=1$、$J=2$ 时超精细劈裂的
示意。量子数 $F$ 可取 $|J-I|,\ldots,J+I-1,J+I$；每个态的能量为
$\frac{A}{2}[F(F+1)-I(I+1)-J(J+1)]$。每个以 $F$ 标记的超精细态有 $2F+1$ 简并，
加磁场 $\mathbf{B}$ 后可按 $m_F$ 分裂。注意朗德间隔定则成立。磁场中劈裂的大小依赖
$F$、$I$、$J$，经由量子数 $F$ 的 $g$ 因子，近似为
$g_F=g_J[F(F+1)+J(J+1)-I(I+1)]/2F(F+1)$。}
\end{figure}

原子（核与电子的组合体）的总角动量 $\mathbf{F}$ 为
\begin{equation*}
\mathbf{F} = \mathbf{J} + \mathbf{I}.\tag{2.70}
\end{equation*}
量子数 $F$ 可取 $|J-I|,\ldots,J+I-1,J+I$。下面几行与自旋--轨道相互作用的处理
几乎完全平行。用式 2.70 可证
\begin{equation*}
\mathbf{F}^{2} = \mathbf{J}^{2}+\mathbf{I}^{2}+2\mathbf{I}\cdot\mathbf{J},\tag{2.71}
\end{equation*}
超精细相互作用取 $A\mathbf{I}\cdot\mathbf{J}$ 形式，其能量期望值为
\begin{equation*}
\langle A\mathbf{I}\cdot\mathbf{J}\rangle = \frac{A}{2}[F(F+1)-I(I+1)-J(J+1)].\tag{2.72}
\end{equation*}
超精细相互作用可作为微扰加入，它把电子能级分裂成以 F 标记的不同超精细结构能级；
每个能级有 $2F+1$ 简并，加磁场可按 $m_F$ 分裂。与式 2.53 直接类比，不同超精细
结构能级的分裂也遵循朗德间隔定则：相邻能级 $E(F)$ 与 $E(F-1)$ 的间隔为
\begin{equation*}
E(F)-E(F-1) = AF.\tag{2.73}
\end{equation*}
即分裂正比于 F——所考虑的相邻两能级中较大的量子数。下面的例子将演示这一点。

\begin{example}{2.5}
图 2.17 对两个例子演示了这一效应。基态原子氢含一个质子（$I=\frac{1}{2}$）和一个
电子（$J=S=\frac{1}{2}$），故 F=0 或 1（图 2.17(a)）。这两态间的能量劈裂极小，
比典型电子能量小约 6 个量级，对应约 1420 MHz 的电磁频率，或约 21 cm 的波长。
这一跃迁是原子氢微波激射器（maser）的基础，在天文学中也很重要：天文学家称之为
H I 发射。当宇宙中原子氢的浓度不足以大量形成 $\mathrm{H_2}$ 分子时，21 cm 辐射
可以热激发产生（所需能量低于宇宙微波背景的温度，见图 3.8）。若气体相对地球运动，
该辐射经历多普勒频移，于是可以由地面测量遥远星系不同区域内气体的转动；它也可用于
我们自己的银河系。

图 2.17(b) 给出更复杂的例子：I=1、J=2 的原子。图中可见朗德间隔定则成立——相邻
能级的劈裂大小随上能级的 F 值按比例变化。
\end{example}

\section*{延伸阅读}

\begin{itemize}[itemsep=0.2em]
\item B.~H.~Bransden 与 C.~J.~Joachain，《Physics of atoms and molecules》，
Longman 1983，提供了关于孤立原子的大量信息。
\item 有用的背景知识还可见 P.~W.~Atkins，《Molecular quantum mechanics》，
OUP 1983。
\item A.~Abragam 与 B.~Bleaney 的全面专著《Electron paramagnetic resonance of
transition ions》，Dover 1986，也很有用。
\item D.~J.~Griffiths，《Introduction to electromagnetism》，Prentice Hall 1989，
对物质中的静磁场给出了易读的论述。
\item 关于抗磁性的经典推导与量子推导优劣的讨论，见 S.~L.~O'Dell 与 R.~K.~P.~Zia，
American Journal of Physics \textbf{54}, 32 (1986)。
\end{itemize}

\section*{习题}

\exer{2.1} 计算基态氢原子气体（数密度 $10^{20}$ m$^{-3}$）的抗磁轨道磁化率，并与
100 K 下的顺磁自旋磁化率比较。

\exer{2.2} 估算一只鸭子的抗磁磁化率（假定它完全由水组成）。要多强的磁场才能在鸭子
中诱导出与一枚磁化铁屑相同的磁矩？对牛重做此题。

\exer{2.3} 计算一块 $\mathrm{CuSO_4\cdot5H_2O}$ 晶体（尺寸
$2\,\mathrm{mm}\times2\,\mathrm{mm}\times2\,\mathrm{mm}$，见表 2.1；密度
$2286\ \mathrm{kg\,m^{-3}}$，相对分子质量 249.7 g）在 10 K、1 T 磁场中的顺磁磁矩。

\exer{2.4} 利用式 2.29 与 2.30 中自旋 $\frac{1}{2}$ 粒子在磁场 B 中的配分函数与
亥姆霍兹函数表达式，证明对单位体积 $n$ 个此类无相互作用粒子，单位体积能量为
\begin{equation*}
E = -n\mu_{\mathrm{B}}B\tanh\left(\frac{\mu_{\mathrm{B}}B}{k_{\mathrm{B}}T}\right),\tag{2.74}
\end{equation*}
单位体积热容为
\begin{equation*}
C = nk_{\mathrm{B}}\left(\frac{\mu_{\mathrm{B}}B}{k_{\mathrm{B}}T}\right)^{2}\operatorname{sech}^{2}\left(\frac{\mu_{\mathrm{B}}B}{k_{\mathrm{B}}T}\right),\tag{2.75}
\end{equation*}
单位体积熵为
\begin{equation*}
S = nk_{\mathrm{B}}\left[\ln\left(2\cosh\left(\frac{\mu_{\mathrm{B}}B}{k_{\mathrm{B}}T}\right)\right) - \frac{\mu_{\mathrm{B}}B}{k_{\mathrm{B}}T}\tanh\left(\frac{\mu_{\mathrm{B}}B}{k_{\mathrm{B}}T}\right)\right].\tag{2.76}
\end{equation*}
这些结果绘于图 2.8。

\exer{2.5} 把上题结果推广到一般 $J$。写配分函数
\begin{equation*}
Z_J(y) = \frac{\sinh[(2J+1)\frac{y}{2J}]}{\sinh[\frac{y}{2J}]}\tag{2.77}
\end{equation*}
其中 $y=xJ=g_J\mu_{\mathrm{B}}JB/k_{\mathrm{B}}T$（见式 2.37），证明熵 $S(y)$ 与
恒场热容 $C(y)$ 作为 $y$ 的函数为
\begin{equation*}
S(y) = nk_{\mathrm{B}}[\ln Z_J(y) - y\,\mathrm{B}_J(y)]\tag{2.78}
\end{equation*}
\begin{equation*}
C(y) = nk_{\mathrm{B}}y^{2}\frac{\mathrm{dB}_J(y)}{\mathrm{d}y},\tag{2.79}
\end{equation*}
$\mathrm{B}_J(y)$ 是布里渊函数。进而证明 $y\ll1$ 时
\begin{equation*}
Z_J(y) = 2J+1\tag{2.80}
\end{equation*}
\begin{equation*}
S(y) = nk_{\mathrm{B}}\left[\ln(2J+1) - \frac{J+1}{3J}y^{2}\right]\tag{2.81}
\end{equation*}
\begin{equation*}
C(y) = nk_{\mathrm{B}}\frac{J+1}{3J}y^{2}\tag{2.82}
\end{equation*}
而 $y\to\infty$ 时
\begin{equation*}
S(y)\to 0\tag{2.83}
\end{equation*}
\begin{equation*}
C(y)\to 0.\tag{2.84}
\end{equation*}

\exer{2.6} 证明：对角动量为 $l$、含 $n$ 个电子的壳层，洪特规则可概括为
\begin{equation*}
\begin{aligned}
S &= \frac{2l+1-|2l+1-n|}{2}\\
L &= S\,|2l+1-n|\\
J &= S\,|2l-n|.
\end{aligned}
\end{equation*}

\exer{2.7} 求下列离子基态的谱项符号：（a）$\mathrm{Ho^{3+}}$（$4f^{10}$）；
（b）$\mathrm{Er^{3+}}$（$4f^{11}$）；（c）$\mathrm{Tm^{3+}}$（$4f^{12}$）；
（d）$\mathrm{Lu^{3+}}$（$4f^{14}$）。

\exer{2.8} 位于原点的点偶极子 $\boldsymbol{\mu}$ 在位置 $\mathbf{r}$ 处产生的磁矢势为
\begin{equation*}
\mathbf{A} = \frac{\mu_0}{4\pi r^{3}}\boldsymbol{\mu}\times\mathbf{r} = \frac{\mu_0}{4\pi}\nabla\times\left(\frac{\boldsymbol{\mu}}{r}\right).\tag{2.85}
\end{equation*}
把它用于原子中电子的哈密顿量
\begin{equation*}
\hat{\mathcal{H}} = \frac{1}{2m_{\mathrm{e}}}(\mathbf{p}+e\mathbf{A})^{2} + 2\mu_{\mathrm{B}}\mathbf{S}\cdot\mathbf{B} + V(r)\tag{2.86}
\end{equation*}
并忽略 $A^{2}$ 项，证明：核偶极子 $\boldsymbol{\mu}=g_I\mu_{\mathrm{N}}\mathbf{I}$
的存在对哈密顿量 $\hat{\mathcal{H}}_0=p^{2}/2m_{\mathrm{e}}+V(r)$ 的微扰
$\hat{\mathcal{H}}'=\hat{\mathcal{H}}-\hat{\mathcal{H}}_0$ 为
\begin{equation*}
\hat{\mathcal{H}}' = \frac{\mu_0}{4\pi}2g_I\mu_{\mathrm{B}}\mu_{\mathrm{N}}\left[(\mathbf{S}\cdot\nabla)(\mathbf{I}\cdot\nabla)\frac{1}{r} - (\mathbf{S}\cdot\mathbf{I})\nabla^{2}\frac{1}{r} + \frac{1}{r^{3}}\mathbf{L}\cdot\mathbf{I}\right].\tag{2.87}
\end{equation*}
证明时需要矢量恒等式
\begin{equation*}
\nabla\times(\nabla\times\mathbf{Q}) = \nabla(\nabla\cdot\mathbf{Q}) - \nabla^{2}\mathbf{Q}\tag{2.88}
\end{equation*}
$\mathbf{Q}$ 为矢量场。利用 $\nabla^{2}(1/r)=-4\pi\delta(\mathbf{r})$（$\delta(\mathbf{r})$
为 delta 函数），证明
\begin{equation*}
(\mathbf{S}\cdot\nabla)(\mathbf{I}\cdot\nabla)\frac{1}{r} = -\frac{1}{r^{3}}\left[\mathbf{S}\cdot\mathbf{I} - \frac{3(\mathbf{S}\cdot\mathbf{r})(\mathbf{I}\cdot\mathbf{r})}{r^{2}}\right] - \frac{4\pi}{3}\mathbf{S}\cdot\mathbf{I}\,\delta(\mathbf{r})\tag{2.89}
\end{equation*}
并由此推出
\begin{equation*}
\hat{\mathcal{H}}' = \frac{\mu_0}{4\pi}2g_I\mu_{\mathrm{B}}\mu_{\mathrm{N}}\mathbf{I}\cdot\left[\frac{\mathbf{L}}{r^{3}} - \frac{\mathbf{S}}{r^{3}} + \frac{3\mathbf{r}(\mathbf{S}\cdot\mathbf{r})}{r^{5}} + \frac{8\pi}{3}\mathbf{S}\,\delta(\mathbf{r})\right].\tag{2.90}
\end{equation*}
方括号中第一项是电子轨道磁矩与核磁矩的耦合；第二、三项是电子自旋与核磁矩的偶极
耦合；最后一项等于费米接触能（式 2.69）。

\exer{2.9} 铂的磁化率为 $2.61\times10^{-4}$，密度 $21450\ \mathrm{kg\,m^{-3}}$，
相对原子质量 $195.09\ \mathrm{g\,mol^{-1}}$。计算其摩尔磁化率（m$^3$\,mol$^{-1}$）
与质量磁化率（m$^3$\,kg$^{-1}$）。用附录 A 把这些结果换算成 cgs 单位：摩尔磁化率
（emu\,mol$^{-1}$）与质量磁化率（emu\,g$^{-1}$）。要理解金属铂的磁化率为何不随
温度变化（与居里定律相反），见第七章。

\exer{2.10} 考虑由式 2.32 的配分函数描述的一组自旋。亥姆霍兹自由能密度为
$F=U-TS=-nk_{\mathrm{B}}T\log Z$，其中 $U=-\mathbf{M}\cdot\mathbf{B}$ 是势能密度，
$S$ 是熵密度。整理可得
\begin{equation*}
TS = -\mathbf{M}\cdot\mathbf{B} + nk_{\mathrm{B}}T\log Z.\tag{2.91}
\end{equation*}
证明：温度不变而磁场改变 $\delta\mathbf{B}$ 时，
\begin{equation*}
T\,\delta S = -\mathbf{B}\cdot\delta\mathbf{M}\tag{2.92}
\end{equation*}
从而
\begin{equation*}
\delta F = \delta U - T\,\delta S = -\mathbf{M}\cdot\delta\mathbf{B}.\tag{2.93}
\end{equation*}
